% Measured inference efficiency supplement; no training or manuscript overwrite.
\paragraph{Inference efficiency.}
We measured the existing best checkpoints from seed 20261001 on an
NVIDIA GeForce RTX 3060 Laptop GPU (6GB), using batch size one, FP32
inputs and weights, no autocast, and evaluation/inference mode.
Matmul TF32 was disabled and cuDNN TF32 was permitted, matching the
existing runtime defaults. Each model was warmed up for 30 calls and
evaluated three times on the same 124 preloaded Validation inputs.
The complete verification pipeline includes the frozen payload decoder
and ViT, evidence fusion, the graph or control module, both heads,
localization upsampling and thresholding. Wall time additionally
includes CPU--GPU input/output transfers and synchronization; file I/O,
resizing, input protection and model loading are excluded.
Peak allocated memory is measured by the PyTorch CUDA allocator and
includes model weights and intermediate tensors, but excludes driver
contexts and other applications. These hardware-specific repeated-input
timings are engineering measurements, not independent-test confidence
intervals or FLOP estimates.

\begin{table*}[!tbp]
\centering
\caption{Inference efficiency on the RTX 3060 Laptop GPU. Mean wall time
is shown with sample SD over 372 timed calls; GPU time is the mean
CUDA-event interval of the network plus GPU postprocessing. Memory
is peak allocated PyTorch CUDA memory (MiB). Checkpoints use the same
training seed; timing repetitions are not training-seed replications.}
\label{tab:e10_inference_efficiency}
\small
\begin{tabular}{@{}lrrrrr@{}}
\toprule
Model & Trainable/M & Total/M & Wall/ms & GPU/ms & Peak MiB \\
\midrule
FULL & 1.416 & 87.242 & $15.36\pm1.66$ & 15.02 & 369.1 \\
NO\_GRAPH & 0.556 & 86.381 & $12.30\pm0.59$ & 11.94 & 360.4 \\
MATCHED\_NO\_GRAPH & 1.416 & 87.241 & $12.79\pm1.24$ & 12.40 & 363.7 \\
SPATIAL\_MIXER\_MATCHED & 1.416 & 87.241 & $13.06\pm1.72$ & 12.68 & 363.7 \\
\bottomrule
\end{tabular}
\end{table*}
